\documentclass[10pt]{article}
\usepackage[T1]{fontenc}
\usepackage{lmodern,microtype,amsmath,amssymb,array,tabularx,booktabs}
\usepackage[margin=0.85in]{geometry}
\usepackage{xurl,hyperref}
\hypersetup{colorlinks=true,urlcolor=blue,pdftitle={Appendix A review -- Version 15}}
\setlength{\parindent}{0pt}
\setlength{\parskip}{5pt}
\newcolumntype{Y}{>{\raggedright\arraybackslash}X}
\newcommand{\Two}{\mathbf2}
\begin{document}
{\LARGE Appendix A: revision and review}\par
{\large Version 15, 31 October 2026}\par

Version 15 makes the Boolean-algebra argument explicit and gives every
downloadable artifact a versioned filename. The standalone appendix is
Appendix A; its mathematical body is Appendix B of the assembled manuscript.
Both PDFs identify version 15 in their metadata and displayed text.

\section*{The categorical construction}
The hypotheses in A.1 include finite products and disjoint, pullback-stable
finite coproducts. Put $\Two=1+1$, with injections
$\mathsf f,\mathsf t:1\to\Two$. Pulling this coproduct back along
$\pi_1:\Two^2\to\Two$ decomposes $\Two^2$ into two copies of
$\Two$. Decomposing each copy again gives
\[
\gamma=[\langle\mathsf f,\mathsf f\rangle,
\langle\mathsf f,\mathsf t\rangle,
\langle\mathsf t,\mathsf f\rangle,
\langle\mathsf t,\mathsf t\rangle]:
1+1+1+1\xrightarrow{\cong}\Two^2.
\]
The chosen copairings define
\[
\neg_2=[\mathsf t,\mathsf f],\qquad
\wedge_2=[\mathsf f,\mathsf f,\mathsf f,\mathsf t]\gamma^{-1},
\qquad
\vee_2=[\mathsf f,\mathsf t,\mathsf t,\mathsf t]\gamma^{-1}.
\]
The coproduct universal property determines a morphism once its restrictions
have been specified. It does not choose those restrictions. The product
decomposition is also essential: it supplies the source needed for a binary
operation on $\Two$.

These morphisms satisfy the Boolean algebra identities. For each finite $n$,
distributivity expresses $\Two^n$ as a coproduct of $2^n$ copies of $1$.
Its summand injections are jointly epimorphic, so equality of two morphisms
from $\Two^n$ can be checked on these injections. Composing the displayed
copairings verifies each axiom. For example,
\[
\wedge_2\langle1_{\Two},\neg_2\rangle=\mathsf f!_{\Two},
\qquad
\vee_2\langle1_{\Two},\neg_2\rangle=\mathsf t!_{\Two}:
\Two\to\Two.
\]
Precomposition with either $\mathsf f$ or $\mathsf t$ proves each equality.
This argument uses the particular coproduct decompositions, not an assumption
that global elements detect arbitrary morphisms. It makes no assertion that
$\Two$ classifies all subobjects.

\section*{Review scope and findings}
A separate agent checked the new A.1 paragraph, the pullback decomposition,
the ordered injections defining $\gamma$, the jointly epimorphic families,
and the types of both complement equations. It found no additional hypothesis
or correction needed for the equational notion of Boolean algebra object.

The four-agent review of version 14 remains a separate record: hostile
mathematical review, mathematical-writing review, six-section outline, and
reconciliation. Those reports are not presented as fresh full reviews of
version 15. Their two principal corrections---explicit copairings and the
initial-algebra composite in the substitution proof---are retained.
The new paragraph explains the mathematical step challenged after that review.

\newpage
\section*{Six-section outline}
\renewcommand{\arraystretch}{1.15}
\begin{tabularx}{\textwidth}{@{}p{0.08\textwidth}>{\raggedright\arraybackslash}p{0.32\textwidth}Y@{}}
\toprule
Section & Title & Function and dependencies\\\midrule
A.1 & Categorical hypotheses and notation & Ambient structure, Boolean
operations and their laws, variable equality, signatures, and parameter stability.\\
A.2 & Terms, formulas, and free variables & Initial syntax and list algebras,
free-variable morphisms, and the sentence subobject; uses A.1.\\
A.3 & Admissible substitution & Single-variable substitution, its admissibility
subobject, and substitution laws; uses A.1--A.2.\\
A.4 & Inference rules and deductively closed subobjects & Inference endofunctor
and the least closed image theorem; needs regularity and finite coproducts.\\
A.5 & First-order theories & Logical schemes, formula and sentence theories,
and partial identities; uses A.2--A.4.\\
A.6 & Propositional theories & The propositional instance of A.4 under weaker
hypotheses, reusing the three propositional schemes of A.5.\\\bottomrule
\end{tabularx}

The section order, theorem labels, bibliography, and dependencies are unchanged.
Numerical realization proofs remain in the applications. The appendix constructs
the theory generated by an axiom subobject under its stated assumptions.

\section*{Publication and verification}
The deliverables have distinct names ending in \texttt{\_v15}: the standalone
PDF and LaTeX source, this review PDF and its source, the assembled manuscript
PDF, and its source archive. The current links point to these versioned files.
The standalone PDF also displays its version immediately below the title.
Earlier repository versions remain available.

The version 14 source retrieved from both GitHub and the saved file already
contained the corrected copairings. Reusing the standalone download path left
the displayed revision ambiguous. Versioned paths and visible PDF version
labels address that publication error.

The final sources were compiled, checked for unresolved references and layout
warnings, and visually inspected on the revised pages. The standalone and
integrated mathematical bodies agree after excluding their publication metadata
and the integrated compatibility label. The accepted abstract and other
mathematical sections are unchanged.

The targeted check is a paper-proof audit. No new full-manuscript, novelty,
bibliography, or Lean verification is claimed. The complete version 14 reports
remain at
\href{https://github.com/flengyel/CategoricalEssentialIncompleteness/tree/main/Audit/2026-10-31-appendix-v14}{the version 14 audit directory}.
\end{document}
